\chapter{Introduction}
\label{chap:introduction}

\section{The Causal-Estimator Evaluation Problem}
\label{sec:introduction:causal-estimator-evaluation-problem}

Causal inference concerns outcomes under alternative actions or exposure conditions, whereas each individual observational record reveals only the factual outcome under the condition that occurred.  The counterfactual outcome needed to score an individual effect estimate is therefore unobserved.  Unlike ordinary prediction, for which a held-out factual outcome can provide a direct target, observational data do not supply a direct causal answer key for evaluating an estimator.  They also leave the true exposure-assignment process and the extent of unmeasured confounding unknown.

Evaluation consequently spans settings with different balances of control and observational fidelity.  Fully synthetic studies specify covariates, exposure assignment, potential outcomes, confounding, support, and effect heterogeneity.  Semi-synthetic studies retain selected empirical components while simulating assignments, outcomes, or both.  Experimentally anchored and empirically anchored designs preserve additional observed structure while obtaining an answer through a controlled experiment, sampling rule, comparison construction, or generated mechanism.  These settings permit causal estimation error to be measured against a specified target, but the evaluator also determines important aspects of the difficulty.  Greater control makes a causal answer measurable, while greater observational fidelity makes the causal target less directly knowable.  These are complementary evaluation properties rather than a simple hierarchy.

CliniCause addresses the observational end of this spectrum.  It constructs resources for examining causal-estimator behavior under source-observed measurement, missingness, prevalence, support, care-process, cohort, and outcome patterns.  Because those resources do not reveal the counterfactual truth, they cannot score causal estimation error against a known effect.  Their role is instead to expose behavior and fragility under realistic observational conditions through several non-equivalent forms of evidence.  Chapter~\ref{chap:background-related-work} situates this role relative to controlled causal benchmarks.

\section{Observational Fidelity in Irregular ICU Records}
\label{sec:introduction:motivation-irregular-icu-time-series-analysis}

Intensive-care medicine produces a particularly demanding form of longitudinal data.  Electronic health records contain laboratory values, vital signs, treatments, administrative fields, and outcomes, but these quantities are neither observed on a common clock nor collected with a common purpose.  Measurements occur at nonuniform times; variables are sampled at markedly different frequencies; many values are absent at most time points; and the intensity of measurement may itself reflect clinical concern, workflow, or resource use.  The resulting records are high dimensional and heterogeneous across patients, care trajectories, and institutions.  Irregular medical time-series methods therefore face more than a standard sequence-modeling inconvenience: they must represent the timing, absence, and changing availability of observations while preserving a defensible analytical link to the question being asked \cite{sun_2026_review_irregular_medical_timeseries,lipton_kale_wetzel_2016_missingness_rnns}.

The conventional convenience of a fixed, complete measurement matrix can conceal these properties.  Regular resampling, imputation, or summary aggregation may be useful engineering choices, but each makes choices about temporal resolution, which observations count, and how absences are represented.  A bedside monitor, a laboratory order, and a static admission field need not convey the same temporal or clinical information.  Moreover, the care process can affect both the values that are available and the chance that they are recorded.  A representation that ignores timing or observation patterns may discard clinically relevant structure; a representation that exploits them can improve prediction while also encoding care-process signals.  This tension makes irregular ICU data a methodological problem spanning statistical representation, workflow design, and interpretation rather than solely a model-selection problem.

For ICU analysis, missingness and irregularity matter at every layer of a study.  They influence which variables a representation can retain, which patterns a prediction model can exploit, which records satisfy a cohort definition, and whether a rule can assign an analytical label.  They also complicate causal adjustment: an absent measurement can represent no measurement, a measurement outside an available window, a source-data limitation, or a clinically informative decision not to measure.  Handling masks, time gaps, or sparse events can improve a predictive representation, but it does not by itself remove the interpretive consequences of the observation process.  The thesis consequently treats irregularity and informative measurement as motivations for careful representation and explicit limitations, rather than as local causal findings.

Two complementary needs follow.  First, an analysis must accommodate irregular longitudinal measurements rather than force all observations into an unqualified regular grid.  Architectures for sparse clinical series and missingness-aware recurrent models illustrate different ways to use values, time intervals, and observation patterns \cite{tipirneni2022strats,che2018grud}.  Second, the analysis needs interpretable variables that can connect a rich event stream to later questions about prognosis and adjusted mortality contrasts.  In this thesis, that intermediate layer consists of \emph{clinically inspired proxy states}: rule-derived analytical constructs that summarize evidence for selected clinical patterns from available measurements.

This wording is deliberate.  A rule-derived proxy label is not a chart-adjudicated diagnosis, ground truth, or a validated phenotype.  A positive tag records that the project-specific source rules were satisfied for the available data; a negative tag does not establish absence of a clinical condition.  Nevertheless, a transparent proxy-state layer can make the passage from irregular measurements to a structured downstream analysis inspectable.  It supplies common targets for multi-label prediction and a documented set of proxy-state exposures for observational analyses, while retaining a visible link to the measurements and rules from which each construct was derived.  Rule-based electronic phenotyping and weak-supervision research provide useful context for this design choice, but do not validate the local proxy definitions or their clinical meaning \cite{banda_2018_electronic_phenotyping,ratner_et_al_2020_snorkel}.

Prediction and causal analysis then answer different questions.  Predicting a rule-derived proxy label asks whether an irregular-time-series model can reproduce that label from the available inputs.  Estimating an adjusted mortality contrast for a proxy-state exposure asks a separate observational question about modeled differences in outcome under explicit assumptions.  Strong performance on the first task does not establish the clinical validity of the label, causal validity of the second analysis, or actionability of the exposure.  The latter task requires a stated causal model, project-specified directed acyclic graphs (DAGs), exposure-specific adjustment logic, attention to support and overlap, sensitivity and permutation diagnostics, and traceable evidence.  Even with these elements, adjusted estimates remain conditional on the adequacy of the proxy measurement, graph, temporal ordering, observed covariates, and model specification \cite{pearl_1995_causal_diagrams,hernan_robins_2016_target_trial,hernan_taubman_2008_well_defined_interventions}.

The framework is instantiated separately on PhysioNet 2012 and MIMIC-III, two public ICU time-series resources that provide complementary settings for this problem \cite{silva2012physionet,johnson2016mimiciii}.  Both contain longitudinal ICU measurements and in-hospital mortality information, yet they differ in era, scale, available variables, data contracts, measurement processes, proxy-state definitions, and project graphs.  These differences make cross-dataset evaluation useful for examining methodological behavior across distinct observational settings, but they preclude treating similarly named constructs as automatically equivalent or pooling numerical estimates without further justification.  The thesis therefore constructs two dataset-specific testbeds through a shared methodological sequence while preserving source-specific definitions and interpretation boundaries.

This problem setting also explains why the thesis moves from data contracts to proxy states, prediction, and only then to adjusted analysis.  The same apparent clinical pattern may have different observable correlates when variable availability, sampling frequency, or coding conventions change.  If a downstream analysis hides those transformations, a reader cannot determine whether a difference arose from the source data, the rule-derived label, the prediction export, or the effect-estimation model.  Making these handoffs explicit is an engineering discipline with scientific consequences: it permits the analysis to be audited, identifies where limitations enter, and prevents an apparently seamless output from being interpreted as a seamless clinical fact.

\section{Representation-Centered Observational Testbeds}
\label{sec:introduction:representation-centered-observational-testbeds}

The central move in CliniCause is to shift controlled researcher intervention away from simulating the patient-level data-generating process and toward specifying an inspectable \emph{representation layer}.  That layer comprises clinically motivated proxy constructs, deterministic operational rules and temporal cutoffs, analytical cohorts, covariates and outcomes, project-specified DAGs, exposure-specific adjustment assumptions, and provenance.  Patient records, measurement and missingness patterns, empirical exposure support, care-process traces, and in-hospital mortality remain source-observed.  CliniCause deterministically instantiates proxy labels from the selected rules; it does not simulate patient records, treatment or exposure-assignment mechanisms, or mortality outcomes.  ``Source-observed'' does not imply that these data are error-free or causally sufficient.

The design and instantiation roles are deliberately separated.  The LLM received schema-level information---available column names and the mortality outcome---together with instructions to use relevant literature.  It did not receive patient records, empirical distributions, exposure prevalences, estimated effects, or execution access.  Its output consisted of literature-grounded design proposals for proxy constructs, operational rules, missingness considerations, and dataset-specific graphs.  Project-selected rules and graphs were subsequently encoded as deterministic source.  The LLM was therefore a design assistant, not a clinical or causal expert, an executed patient-level component, or a source of ground truth.

This representation is instantiated independently on MIMIC-III and PhysioNet 2012.  The resulting resources have a common analytical role but are not pooled replicas of one object.  Similarly named proxy constructs can differ in measurement availability, operational definition, prevalence, missingness, support, graph structure, adjustment, and population.  Each resource is a \emph{testbed} when used to study causal-analysis behavior; the term does not imply a benchmark with known causal truth.

Figure~\ref{fig:clinicause-testbed-pipeline} summarizes the separation among design, deterministic instantiation, and evaluation.

\begin{figure}[htbp]
  \centering
  \includegraphics[width=\textwidth]{clinicause_testbed_pipeline.png}
  \caption{The Design, Instantiate, and Evaluate stages of CliniCause.  Schema and literature information support LLM-assisted representation proposals, which project-selected deterministic code instantiates at patient level on source-observed ICU records and outcomes.  The resulting dataset-specific resources are characterized through converging evidence without a known causal answer key; the LLM receives schema information rather than patient records.}
  \label{fig:clinicause-testbed-pipeline}
\end{figure}

\section{Thesis Gap and Objective}
\label{sec:introduction:thesis-gap-objective}

Prior work offers important components for this setting.  Irregular-time-series research provides representations and predictive methods for sparse, asynchronously observed medical data.  Electronic phenotyping provides strategies for constructing analytical variables when direct expert adjudication is unavailable \cite{banda_2018_electronic_phenotyping}.  Data programming formalizes programmatic heuristic labeling \cite{ratner_et_al_2016_data_programming}, while Snorkel provides related weak-supervision and downstream discriminative-training context \cite{ratner_et_al_2020_snorkel}.  Evaluations of medical LLM capabilities provide a bounded precedent for design-time clinical-knowledge elicitation \cite{singhal_et_al_2023_llm_clinical_knowledge}; work on LLM-derived causal-graph priors separately motivates their use as proposals rather than validated graphs \cite{darvariu_et_al_2024_llm_causal_graph_priors}.  Causal-inference research supplies DAG-based reasoning, double machine learning, heterogeneous-effect estimation, and diagnostics for observational studies \cite{chernozhukov2018dml,wager2018causalforest,athey2019grf}.  ICU-focused reviews further emphasize the difficulty of defining an actionable question, managing confounding, and distinguishing predictive patterns from treatment-effect claims in observational clinical records \cite{smit_2023_causal_inference_icu_scoping_review,bica_2021_individualized_treatment_effects_ehr_ml}.

The gap addressed here is not merely the absence of an integrated pipeline.  Controlled benchmarks provide measurable causal accuracy under specified mechanisms, while routine observational data provide realistic processes without a known causal answer.  A reusable observational evaluation object must make the intervening representation decisions explicit: LLM-generated proposals must be distinguished from selected source-encoded designs; source-specific preprocessing must yield well-defined records; rule-derived labels must be distinguishable from model predictions and algorithmic aggregates; prediction exports must be normalized before downstream use; and causal analyses must state how their exposures, covariates, graphs, diagnostics, populations, and provenance relate.  Without these interfaces, a resource can appear coherent while its design assumptions, labels, analysis rows, adjustment variables, or interpretation change silently between stages.

The objective of this thesis is therefore \emph{to construct and characterize clinically grounded observational causal-analysis testbeds from source-observed irregular ICU records through an evidence-tracked, representation-centered framework, while making their proxy definitions, data contracts, causal assumptions, analytical interfaces, diagnostics, provenance, and limitations explicit}.  ``Evidence-tracked'' here means that the thesis distinguishes LLM design proposals, selected source-coded behavior, checked numerical artifacts, and unresolved provenance or validation requirements.  The testbeds and their enabling methodology constitute an observational evaluation contribution; they are not claims of clinical validation, deployment readiness, causal identification, or full clean-checkout computational reproducibility.

The framework has two distinct layers.  At design time, a structured LLM-assisted protocol elicited proposed proxy-state ontologies, decision-rule families, missingness considerations, and dataset-specific DAGs; project-selected designs were then encoded in deterministic tagger and graph source files.  During analytical execution, the sequence is dataset-specific preprocessing; deterministic generation of rule-derived proxy labels; multi-label proxy-state prediction; prediction normalization and algorithmic aggregation; exposure-specific adjustment under the project-specified DAGs; matching and heterogeneous-effect estimation; and robustness and provenance evaluation.  The LLM is not invoked in this executed patient-level pipeline.  The proxy states are not validated clinical diagnoses; the DAGs are not learned from data; no unconditional causal effect is claimed; and many illness-state proxies do not define well-specified interventions.  The workflow does not provide clinical recommendations, bedside decision support, or prospective evidence of benefit.

The integrated workflow is the enabling machinery that makes the testbeds inspectable rather than assumed.  Dataset-specific preprocessing preserves the distinction between source records and analysis-ready objects.  Rule-based construction makes the origin of each proxy label visible.  Prediction normalization distinguishes probabilities from thresholded labels, and algorithmic aggregation distinguishes a deterministic vote from clinical consensus.  Project-specified DAGs and exposure-specific adjustment sets make the causal assumptions operationally visible, while matching and model-estimated CATE summaries provide different forms of observational evidence.  Finally, robustness and provenance checks separate a checked numerical statement from stronger claims about causal identification, clinical meaning, or rerun reproducibility.  The framework's purpose is not to erase these distinctions, but to retain them throughout construction and characterization of the resources.

This framing has practical consequences for how the results should be read.  A checked prediction metric concerns the archived evaluation of a model against rule-derived labels, not an independently verified clinical endpoint.  A matching output is a descriptive matched-pair outcome difference, not automatically a treatment effect.  A mean model-estimated CATE is an aggregate of fitted conditional estimates under the project assumptions, not an unconditional population quantity.  Sensitivity, permutation, and support diagnostics contribute evidence about particular failure modes, but none can substitute for a justified intervention definition, an adequate causal graph, or clinical review.  By placing these distinctions in the framework itself, the thesis avoids treating provenance, statistical modeling, and clinical interpretation as afterthoughts to an otherwise complete pipeline.

LLM-assisted clinical and causal knowledge elicitation is therefore a substantive design-time stage of CliniCause.  The structured protocol produced proposals for the proxy-state ontology, binary decision rules, missingness reasoning, and dataset-specific DAGs; the project selected designs and encoded them in active deterministic source code.  This role does not make the LLM a validated medical expert, establish that every proposal was implemented unchanged, or validate the proxy constructs or graphs.  Nor was the LLM an executed patient-level estimator: it did not preprocess records, generate runtime predictions, select adjustment variables dynamically, or estimate causal effects.  The distinction between design provenance and executable authority makes the representation layer inspectable without treating generated proposals as clinical or causal evidence.

\section{Research Questions and Contributions}
\label{sec:introduction:research-questions-contributions}

The main research question is:

\begin{quote}
How can clinically grounded observational causal-analysis testbeds be constructed from source-observed irregular ICU records---without simulating patient records, treatment or exposure-assignment mechanisms, or outcomes---and what converging evidence can characterize their usefulness and limitations for evaluating causal estimators?
\end{quote}

The words ``observational'' and ``characterize'' are central.  The thesis investigates how a linked workflow can construct and qualify retrospective evaluation resources; it does not claim to score estimates against causal truth, prove that proxy-state exposures cause mortality, or show that manipulating a proxy would improve outcomes.  The main question is addressed through five secondary research questions:

\begin{description}
  \item[SRQ-1:] How can clinically and literature-grounded representation designs be instantiated deterministically and separately in MIMIC-III and PhysioNet while retaining source-specific measurement and missingness processes?
  \item[SRQ-2:] Which data contracts, normalization, aggregation, DAG, adjustment, and provenance mechanisms make the constructed resources explicit, auditable, and reusable?
  \item[SRQ-3:] How recoverable are the deterministic proxy labels from irregular time series, and what mortality-relevant predictive information is retained by the constructed proxy representations?
  \item[SRQ-4:] What converging evidence---including estimator sign, rank, and magnitude agreement, clinical-literature comparison, matching and support, permutation disruption, and omitted-variable sensitivity---characterizes each resource?
  \item[SRQ-5:] Which limitations prevent the resources from being interpreted as clinical validation, known causal truth, or proof of estimator correctness?
\end{description}

The primary contribution is a representation-centered framework and two concrete, dataset-specific observational causal-analysis testbeds constructed from source-observed MIMIC-III and PhysioNet 2012 records.  The framework preserves observed records and outcomes while making its proxy definitions, operational rules, cohort choices, DAGs, adjustment assumptions, and provenance explicit and replaceable.  Because no known causal answer is available, it characterizes the resources through converging, non-equivalent evidence rather than estimator error against ground truth.

The enabling methods include schema-bounded LLM-assisted elicitation and deterministic source code for proxy rules and DAGs.  They also include dataset-specific data contracts, irregular-time-series proxy prediction, deterministic aggregation, exposure-specific DAG adjustment, matching, heterogeneous-effect estimation, diagnostics, and provenance.  The compared estimators are CausalForestDML, LinearDML, and CausalPFN.  CausalForestDML remains the primary model-estimated causal-analysis estimator, LinearDML is a secondary comparator, and CausalPFN is retained as an exploratory complementary estimator because it lacks the complete diagnostic support available for the DML estimators.

The empirical contribution is a checked cross-dataset comparison of proxy recoverability using four learned prediction models, original-cohort analyses for all 19 dataset--exposure combinations, cross-estimator directional comparison, separate outcome-downsampling robustness analyses, and cross-dataset characterization without numerical pooling.  CausalForestDML and LinearDML agree in direction for 19 of 19 combinations; all three estimators agree for 18 of 19, with PhysioNet shock retained as the exception.  Matching and support diagnostics, mortality-relevant prediction, clinical-literature comparison, permutation disruption, omitted-variable sensitivity, population perturbation, and provenance provide additional evidence classes whose numerical findings and qualifications are reported in their designated chapters.  The thesis also contributes checked numerical tables, a results manifest, checksums, a source packet, a decision register, corrected result figures, and explicit tracking of unresolved provenance.  These materials support numerical traceability, but they do not establish that the analysis can be rerun from a clean checkout or that every historical producing setting has been recovered.

These contributions are cumulative rather than interchangeable.  Construction of the two observational testbeds is primary; the integrated workflow and shared proxy/data-contract layers make their handoffs inspectable; empirical comparisons and estimator triangulation characterize the resources; and evidence records preserve the boundary between numerical traceability and full computational reproducibility.  Agreement can show that a pattern is not unique to one fitted estimator under one representation, but it does not establish estimator accuracy, graph correctness, causal identification, or clinical validity.  None of these layers upgrades a rule-derived label to a diagnosis or a graph-guided contrast to a well-defined intervention effect.

As empirical orientation, the four-model comparison was dataset dependent: STraTS led the four archived MIMIC-III test metrics and GRU-D led the corresponding PhysioNet metrics.  The original-cohort directional agreements summarized above motivate the later analysis, while the PhysioNet shock exception remains visible.  These patterns do not establish statistical superiority, causal validity, positivity, or clinical utility; matching and outcome-downsampled results remain qualified robustness evidence.

\section{Thesis Organization}
\label{sec:introduction:thesis-organization}

The remainder of the thesis follows construction and characterization of the two testbeds from context to interpretation.  Background and study design establish the evaluation spectrum, resource contracts, questions, and assumptions.  Chapters~\ref{chap:data-preprocessing}--\ref{chap:experimental-design} construct and characterize the resources through source-specific contracts, proxy instantiation and prediction, DAG-guided adjustment and estimation, diagnostics, and provenance.  The final chapters report the checked evidence, interpret its limitations, and state the conclusions and future-work priorities.
